\documentclass{beamer}
\usetheme{Madrid}
\usecolortheme{default}

\usepackage{graphicx}
\usepackage{booktabs}

\definecolor{uoftblue}{RGB}{6,41,88}
\setbeamercolor{titlelike}{bg=uoftblue}
\setbeamerfont{title}{series=\bfseries}

\title[Adaptive Chunking for RAG]
{Enhancing RAG with Adaptive Chunking}

\subtitle{Retrieval Benchmark and Lightweight Knowledge Graph Prototype}

\author[Group Project]
{Group Members}

\institute[]
{
  Department of Electrical and Computer Engineering\\
  University of Toronto
}

\date[Winter 2026]
{ECE1508 Applied Deep Learning -- Winter 2026}

\begin{document}

\frame{\titlepage}

\begin{frame}
\frametitle{Problem and Our Approach}
\small
\begin{itemize}
  \item Standard RAG often fails when evidence is split across chunks or linked by latent entity relations.
  \item Our proposal direction: adaptive chunking + stronger retrieval + graph-based reasoning.
  \item Implemented system:
  \begin{itemize}
    \item KILT NaturalQuestions benchmark pipeline
    \item Bi-encoder baseline: MiniLM + FAISS
    \item Late interaction retriever: ColBERTv2
    \item Chunking ablation: paragraph, sentence-window, adaptive-sentence
    \item Lightweight entity knowledge graph prototype
  \end{itemize}
  \item Goal of this milestone: build a truthful, reproducible retrieval foundation before full end-to-end RAG.
\end{itemize}
\end{frame}

\begin{frame}
\frametitle{Benchmark Design}
\small
\begin{itemize}
  \item Dataset: KILT NQ validation, grouped into single-entity and multi-entity queries by spaCy NER.
  \item Reduced corpus: sampled gold Wikipedia pages + distractors.
  \item Metrics: Recall@1/5/10/20, retrieval latency, index size.
  \item Why Recall@k:
  \begin{itemize}
    \item retrieval-first milestone
    \item directly measures whether correct evidence is retrieved
    \item useful before adding generation and RAGAS
  \end{itemize}
  \item Key retriever difference:
  \begin{itemize}
    \item Bi-encoder compresses each chunk into one vector
    \item ColBERT keeps token-level representations and uses late interaction
  \end{itemize}
\end{itemize}
\end{frame}

\begin{frame}
\frametitle{Chunking Ablation Results}
\small
\begin{itemize}
  \item Quick ablation over 51 queries:
  \begin{itemize}
    \item paragraph: 3415 chunks
    \item sentence-window: 4481 chunks
    \item adaptive-sentence: 3511 chunks
  \end{itemize}
  \item Best overall ColBERT performance came from adaptive-sentence:
  \begin{itemize}
    \item Recall@1: 0.1134 $\rightarrow$ 0.1199
    \item Recall@5: 0.2487 $\rightarrow$ 0.2553
    \item Recall@10: 0.2949 $\rightarrow$ 0.2977
  \end{itemize}
  \item Interpretation:
  \begin{itemize}
    \item adaptive chunking helps ColBERT more than the bi-encoder
    \item sentence-window costs more chunks without a clear gain over adaptive-sentence
    \item gains are selective, not universal
  \end{itemize}
\end{itemize}
\end{frame}

\begin{frame}
\frametitle{Knowledge Graph Prototype and Findings}
\small
\begin{itemize}
  \item Built a lightweight graph from chunk-level entity extraction:
  \begin{itemize}
    \item entity $\rightarrow$ chunk links
    \item entity $\leftrightarrow$ entity co-occurrence edges
  \end{itemize}
  \item Graph retrieval is weaker than semantic retrieval overall:
  \begin{itemize}
    \item overall Recall@10 = 0.0878
  \end{itemize}
  \item But the graph baseline is more aligned with multi-entity questions:
  \begin{itemize}
    \item multi-entity Recall@10 = 0.1259
    \item single-entity Recall@10 = 0.0611
  \end{itemize}
  \item Takeaway: the current graph prototype is not yet a replacement for dense retrieval, but it is promising for relation-heavy queries.
\end{itemize}
\end{frame}

\begin{frame}
\frametitle{Conclusions and Next Step}
\small
\begin{itemize}
  \item We now have a full retrieval benchmark pipeline that runs end to end.
  \item Main findings:
  \begin{itemize}
    \item ColBERT improves top-ranked retrieval over the bi-encoder baseline.
    \item Adaptive-sentence chunking is our best current chunking strategy.
    \item The graph prototype is weak overall but more useful for multi-entity queries.
  \end{itemize}
  \item Most valuable next step:
  \begin{itemize}
    \item combine semantic retrieval with graph signals, instead of using the graph alone
  \end{itemize}
  \item Honest project status:
  \begin{itemize}
    \item retrieval benchmark and KG prototype are implemented
    \item end-to-end RAG, Neo4j, routing, and RAGAS remain future work
  \end{itemize}
\end{itemize}
\end{frame}

\end{document}
